%% %%%%%%%%%%%%%%%%%%%%%%%%%%%%%%%%%%%%%%%%%%%%%%%%%%%%%%%%%%%%%%%
%% AGFA-Book-Einleitung.tex -- Kapitel ohne Nummer
%% \addchap steht im Inhaltsverzeichnis und im Kolumnentitel,
%% bekommt aber keine Nummer.
%% %%%%%%%%%%%%%%%%%%%%%%%%%%%%%%%%%%%%%%%%%%%%%%%%%%%%%%%%%%%%%%%
\addchap{Einleitung}
\label{chap:einleitung}

Diese Vorlage zeigt den Aufbau einer Arbeit mit Kapiteln. In
\cref{chap:grundlagen} werden die Begriffe eingeführt, in
\cref{chap:anwendungen} angewendet. Das Hauptresultat ist
\cref{thm:banach}; es stützt sich auf \cref{lem:cauchy}.

Innerhalb eines Kapitels heißen Sätze und Gleichungen kurz
\enquote{Theorem 2.1} bzw.\ \enquote{(2.1)}: Die erste Zahl ist die
Section, die zweite zählt in der Section. Bei Verweisen wird die
Kapitelnummer vorangestellt, also \enquote{Theorem I.2.1}. So ist jeder
Verweis eindeutig, egal von wo er kommt.
